% Options for packages loaded elsewhere
\PassOptionsToPackage{unicode}{hyperref}
\PassOptionsToPackage{hyphens}{url}
%
\documentclass[
  11pt,
]{article}
\usepackage{lmodern}
\usepackage{amssymb,amsmath}
\usepackage{ifxetex,ifluatex}
\ifnum 0\ifxetex 1\fi\ifluatex 1\fi=0 % if pdftex
  \usepackage[T1]{fontenc}
  \usepackage[utf8]{inputenc}
  \usepackage{textcomp} % provide euro and other symbols
\else % if luatex or xetex
  \usepackage{unicode-math}
  \defaultfontfeatures{Scale=MatchLowercase}
  \defaultfontfeatures[\rmfamily]{Ligatures=TeX,Scale=1}
\fi
% Use upquote if available, for straight quotes in verbatim environments
\IfFileExists{upquote.sty}{\usepackage{upquote}}{}
\IfFileExists{microtype.sty}{% use microtype if available
  \usepackage[]{microtype}
  \UseMicrotypeSet[protrusion]{basicmath} % disable protrusion for tt fonts
}{}
\makeatletter
\@ifundefined{KOMAClassName}{% if non-KOMA class
  \IfFileExists{parskip.sty}{%
    \usepackage{parskip}
  }{% else
    \setlength{\parindent}{0pt}
    \setlength{\parskip}{6pt plus 2pt minus 1pt}}
}{% if KOMA class
  \KOMAoptions{parskip=half}}
\makeatother
\usepackage{xcolor}
\IfFileExists{xurl.sty}{\usepackage{xurl}}{} % add URL line breaks if available
\IfFileExists{bookmark.sty}{\usepackage{bookmark}}{\usepackage{hyperref}}
\hypersetup{
  hidelinks,
  pdfcreator={LaTeX via pandoc}}
\urlstyle{same} % disable monospaced font for URLs
\usepackage[margin=1in]{geometry}
\setlength{\emergencystretch}{3em} % prevent overfull lines
\providecommand{\tightlist}{%
  \setlength{\itemsep}{0pt}\setlength{\parskip}{0pt}}
\setcounter{secnumdepth}{-\maxdimen} % remove section numbering
\usepackage{esint}
\setmathfont{Asana Math}
\setlength{\emergencystretch}{3em}

\author{}
\date{}

\begin{document}

\hypertarget{exact-nonlinear-slow-mode-energy-and-fixed-amplitude-adiabatic-metrics}{%
\section{Exact nonlinear slow-mode energy and fixed-amplitude adiabatic
metrics}\label{exact-nonlinear-slow-mode-energy-and-fixed-amplitude-adiabatic-metrics}}

24 September 2026. This note replaces a linearized fiber deformation by
an actual horizontal curve of hyperbolic metrics. It determines every
Ricci block, the complete nonlinear slow equation, and its normal
feedback. The resulting energy identity remains valid at fixed, nonzero
moduli amplitude. On cyclic covers it produces pointwise and
unnormalized integral Ricci errors tending to zero while that amplitude
does not shrink.

The full-mode theorem in
\href{20260924_integrable_kernel_cyclic_repair.md}{I33} concerns
arbitrary tensor perturbations in a shrinking neighborhood and measures
distance to \textbf{all} normalized Einstein metrics. The present
theorem treats a larger nonlinear, geometrically specified family, but
its fixed-distance conclusion concerns fiber-preserving product outputs.
Neither conclusion replaces the other. The underlying smooth products
already carry exact Einstein metrics.

The submersion identities and the exact adapted rigidity statement were
already available in
\href{../mathematics/current/round20_construction.md}{Round 20}. They
are not counted as new results. The additions here are the exact
arclength equation, the full residual-energy identity including normal
acceleration, and its fixed-amplitude application. The companion
\href{20260924_moduli_energy_gauge.md}{actual-gauge identity} places
these metrics inside the first line's actual bilateral monopole
near-saturation regime.

\hypertarget{an-actual-horizontally-marked-hyperbolic-family}{%
\subsection{1. An actual, horizontally marked hyperbolic
family}\label{an-actual-horizontally-marked-hyperbolic-family}}

Let \(F\) be a closed oriented surface of genus at least two. Fix a
smooth curve \(h_v\), defined on an open interval containing a compact
interval \(J\), of metrics of Gaussian curvature \(-1\). Require \[
 k_v=\partial_vh_v,\qquad
 \operatorname{tr}_{h_v}k_v=0,\qquad
 \operatorname{div}_{h_v}k_v=0,\qquad
 \int_F|k_v|_{h_v}^2\,dA_{h_v}=1.
 \tag{1}
\] Such nonconstant curves exist. A regular curve in Teichmüller space
has a horizontal lift; after reparametrization by its real-tensor
\(L^2\) arclength it satisfies (1). Rupflin--Topping, Definition 1.1 and
the discussion of equation (1.1) in
\href{https://link.springer.com/article/10.1007/s00526-018-1386-y}{\emph{Horizontal
curves of hyperbolic metrics}}, give the horizontal decomposition.
Horizontality is equivalently the real part of a holomorphic quadratic
differential. Its length is a fixed normalization of Weil--Petersson
length; no conventional WP constant is silently inserted into (1).

Pointwise trace-freeness implies that all \(h_v\) have the \textbf{same
area form}, not just the same total area. Write this form as \(dA_F\),
and set \[
 V_F=\int_FdA_F=4\pi(\operatorname{genus}(F)-1),\qquad
 N_v=\partial_vk_v-k_vh_v^{-1}k_v,
\] \[
 C_4(v)=\int_F|k_v|_{h_v}^4dA_F,\qquad
 C_N(v)=\int_F|N_v|_{h_v}^2dA_F,\qquad
 \Xi(v)=\frac{C_4(v)}{16}+\frac{C_N(v)}4.
 \tag{2}
\] Here a juxtaposition such as \(k_vh_v^{-1}k_v\) denotes contraction
as bilinear forms. All these are tensors and functions of the actual
marked metrics on \(F\), not only formal Teichmüller coordinates.
Differentiating the trace and unit-speed conditions gives \[
 \operatorname{tr}_{h_v}N_v=0,\qquad
 \int_F\langle N_v,k_v\rangle_{h_v}dA_F=0,
 \qquad \Xi(v)\ge\frac1{16V_F}.
 \tag{3}
\] For example, with \(A=h_v^{-1}k_v\), differentiation of the squared
\(L^2\) norm yields
\(0=2\int_F(\langle\partial_vk_v,k_v\rangle-\operatorname{tr}A^3)\).
This proves the middle identity. The last follows from Cauchy--Schwarz
and \(\int_F|k_v|^2=1\). These identities do not assume that the curve
is a WP geodesic. In particular, its normal acceleration is not
discarded.

Let \((B,\gamma)\) be any closed connected oriented Riemannian surface,
and let \(v:B\to J\) be smooth. On the original smooth product define \[
             G=\gamma+h_{v(x)},\qquad E=\operatorname{Ric}_G+G.
 \tag{4}
\] The marking in (1) defines a flat horizontal distribution.
Throughout, \(\Delta_\gamma=-\operatorname{div}_\gamma\nabla\) is
nonnegative, and tensor norms without an indicated average are
unnormalized norms.

\hypertarget{complete-nonlinear-equation-and-residual-energy}{%
\subsection{2. Complete nonlinear equation and residual
energy}\label{complete-nonlinear-equation-and-residual-energy}}

\textbf{Theorem A.} For (1)--(4), all Ricci-error blocks are exactly \[
 \begin{aligned}
 E_{HH}&=(K_\gamma+1)\gamma
                      -\tfrac14|k_v(y)|_{h_v}^2\,dv\otimes dv,\\
 E_{HV}&=0,\\
 E_{VV}&=\tfrac12 k_v\Delta_\gamma v
                         -\tfrac12N_v|dv|_\gamma^2.
 \end{aligned}
 \tag{5}
\] Consequently the retained slow equation and its entire transverse
component along this family are \[
 \int_F\langle E_{VV},k_v\rangle dA_F
            =\tfrac12\Delta_\gamma v,
 \qquad
 E_{VV}-k_v\int_F\langle E_{VV},k_v\rangle dA_F
            =-\tfrac12N_v|dv|^2.
 \tag{6}
\] They give the exact global identity \[
 \boxed{\begin{aligned}
 \int_{B\times F}|E|_G^2dV_G
 =\int_B\bigg[&\frac14(\Delta_\gamma v)^2
  +2V_F\left(K_\gamma+1-\frac{|dv|^2}{8V_F}\right)^2\\
  &+\left(\Xi(v)-\frac1{32V_F}\right)|dv|^4\bigg]dA_\gamma.
 \end{aligned}}
 \tag{7}
\] In particular, for an arbitrary base metric, \[
 \int|E|^2\ge\frac14\int_B(\Delta_\gamma v)^2dA_\gamma
                  +\frac1{32V_F}\int_B|dv|^4dA_\gamma.
 \tag{8}
\] If \(K_\gamma=-1\), the simpler identity is \[
 \boxed{\quad
 \int|E|^2=\frac14\int_B(\Delta_\gamma v)^2dA_\gamma
                    +\int_B\Xi(v)|dv|^4dA_\gamma.
 \quad}
 \tag{9}
\] There is no potential term depending only on \(v\), and no uncomputed
\(O(v^2)\) error. The exact family of hyperbolic fibers is responsible
for this cancellation.

\textbf{Proof.} In base coordinates \(x^i\) and fiber coordinates
\(y^a\), put
\(S_i=\tfrac12h^{-1}\partial_i h=\tfrac12h^{-1}k_v\partial_i v\). Its
trace is zero. The mixed Christoffel symbols of the block metric are \[
 \Gamma^a_{ib}=\tfrac12(h^{-1}k_v)^a{}_b\partial_i v,
 \qquad
 \Gamma^i_{ab}=-\tfrac12\gamma^{ij}(k_v)_{ab}\partial_jv;
 \tag{10}
\] the remaining nonzero symbols within each factor are its intrinsic
Christoffel symbols. Substitution into Ricci, or the full
flat-connection identities in Round 20, gives \[
 (\operatorname{Ric}_G)_{ij}
    =K_\gamma\gamma_{ij}-\operatorname{tr}(S_iS_j),
 \qquad (\operatorname{Ric}_G)_{ia}
    =(\operatorname{div}_{h_v}S_i)_a,
\] \[
  (E_{VV})^\sharp=-\gamma^{ij}D_iS_j,
  \quad
  D_iS_j=\partial_iS_j+[S_i,S_j]-\Gamma^\ell_{ij}(\gamma)S_\ell.
 \tag{11}
\] The mixed term vanishes by (1). The commutator vanishes because every
\(S_i\) is a scalar multiple of the same endomorphism. Differentiating
\(h^{-1}k_v\) gives \(h^{-1}N_v\). Since
\(\gamma^{ij}(\partial_i\partial_jv-\Gamma^\ell_{ij}\partial_\ell v) =-\Delta_\gamma v\),
(11) proves all of (5).

Equation (3) now proves (6) and \[
 \int_F|E_{VV}|^2dA_F
           =\tfrac14(\Delta_\gamma v)^2
                          +\tfrac14C_N(v)|dv|^4.
 \tag{12}
\] Set \(c=K_\gamma+1\), \(X=|dv|^2\). Since \(|dv\otimes dv|^2=X^2\)
and \(\langle\gamma,dv\otimes dv\rangle=X\), the horizontal norm
integrated over a fiber is \[
 2V_Fc^2-\tfrac12cX+\tfrac1{16}C_4(v)X^2.
 \tag{13}
\] Completing this scalar square and using (12) proves (7). Equations
(8)--(9) follow from (3) and \(c=0\), respectively. \(\square\)

\textbf{Corollary A1 (forced slow energy and exact classification).}
Define \(f(x)=\int_F\langle E_{VV},k_v\rangle dA_F\). Then, for any
constant \(c\), \[
 \Delta_\gamma v=2f,\qquad
 \int_B|dv|^2dA_\gamma=2\int_B(v-c)f\,dA_\gamma,
\] \[
 \|v-\bar v\|_{L^2(B)}
       \le\frac{2}{\lambda_1(B,\gamma)}\|E\|_{L^2(B\times F)}.
 \tag{14}
\] Here \(\lambda_1>0\) is the first nonzero scalar eigenvalue, and the
last inequality follows from \(\|f\|_2\le\|E_{VV}\|_2\). No scalar
spectral gap has been silently made uniform in the base.

If \(G\) is Einstein with \textbf{any} constant \(\lambda\), then
\(\operatorname{tr}_{VV}\operatorname{Ric}_G=-2\) by (3),(5), so
\(\lambda=-1\). Equivalently, pairing the Einstein equation with \(k_v\)
in (6) already implies \(\Delta v=0\). Connectedness of the closed base
gives \(v\) constant, and then (5) forces \(K_\gamma=-1\). Thus the
exact solutions in this entire block family are precisely the products.
This recovers the relevant inherited adapted rigidity, while (6)--(14)
supply its quantitative nonlinear equation and full residual energy.

\textbf{Corollary A2 (local stress and global amplitude).} Put
\(\delta=\|G^{-1}E\|_{\infty,\mathrm{op}}\). For an arbitrary base, \[
 |dv|^2\le8V_F\delta,
 \qquad
 \operatorname{osc}_B v\le\sqrt{8V_F\delta}\,
                                  \operatorname{diam}(B,\gamma).
 \tag{15}
\] When \(K_\gamma=-1\), 8 can be replaced by 4. Indeed, the trace-free
part of the horizontal block has operator norm \(|k_v|^2|dv|^2/8\); its
norm is at most \(\delta\) in dimension two. Taking a fiber maximum and
using \(\int_F|k_v|^2=1\) proves the first bound. For a hyperbolic base
the whole horizontal error has norm \(|k_v|^2|dv|^2/4\). Integrating
along a minimizing base geodesic proves the amplitude estimate.
Base-curvature adjustment therefore cannot remove the rank-one
trace-free stress.

\hypertarget{fixed-amplitude-cyclic-cover-construction}{%
\subsection{3. Fixed-amplitude cyclic-cover
construction}\label{fixed-amplitude-cyclic-cover-construction}}

Fix a closed hyperbolic surface \((B,b)\) and a primitive class in
\(H^1(B;\mathbb Z)\). Let \(\eta\) be its harmonic representative,
normalized to have integral periods. On the associated infinite cyclic
cover \(B_\infty\), write \(dz=\eta\), with \(z(\tau x)=z(x)+1\) for a
deck generator \(\tau\). The function \(z\) is harmonic. Let
\(B_m=B_\infty/\langle\tau^m\rangle\), with the lifted hyperbolic metric
\(b_m\).

Fix \(a>0\) with \([-a,a]\subset J\), and suppose the two marked
hyperbolic structures \([h_{-a}]\) and \([h_a]\) are different. This is
achievable on an arbitrarily short regular horizontal arc. Define the
actual metrics \[
 v_m=a\cos(2\pi z/m),\qquad
 M_m=B_m\times F,\qquad G_m=b_m+h_{v_m}.
 \tag{16}
\] The amplitude \(a\) is fixed independently of \(m\).

\textbf{Theorem B.} This sequence has uniform local smooth bounded
geometry and is noncollapsed on a fixed local scale. For every fixed
integer \(j\ge0\) and \(0<\alpha<1\), in the uniformly lifted product
charts, \[
 \|\operatorname{Ric}_{G_m}+G_m\|_{C^{j,\alpha}_{\mathrm{uloc}}}
                  \le C_{a,j,\alpha}m^{-2},
 \qquad
 c_a m^{-2}\le\delta_-(G_m)\le C_a m^{-2},
 \tag{17}
\] \[
 c_a m^{-3}\le
       \int_{M_m}|\operatorname{Ric}_{G_m}+G_m|^2dV_{G_m}
                 \le C_a m^{-3},\qquad
 \operatorname{osc}_{B_m}v_m=2a.
 \tag{18}
\] The constants depend on the fixed surfaces, harmonic form, and fixed
horizontal curve, not on the covering degree. The volume is exactly
\(m\operatorname{Area}(B,b)V_F\), and diameter is comparable to \(m\).
In particular, even unnormalized \(L^2\) Ricci error tends to zero.

\textbf{Proof.} The harmonic chain rule gives exactly \[
 dv_m=-a\frac{2\pi}{m}\sin(2\pi z/m)\eta,
 \qquad
 \Delta_{b_m}v_m=\left(\frac{2\pi}{m}\right)^2
                                      |\eta|^2v_m.
 \tag{19}
\] All derivatives of the lifted \(\eta\) are uniformly bounded. Every
fixed number of derivatives of \(|dv_m|^2\) and \(\Delta v_m\) is
therefore \(O(m^{-2})\). The compact smooth hyperbolic path has uniform
bounds of every fixed order, so (5) proves the first part of (17).
Although the Hessian of \(v_m\) need only be \(O(m^{-1})\), the complete
Ricci formula involves its \textbf{trace} in (19); there is no missing
\(O(m^{-1})\) mixed block, because horizontality makes it vanish
exactly.

The metrics \(G_m\) are uniformly equivalent to \(b_m+h_0\), with all
local coefficient derivatives uniformly bounded. Lift a fixed finite
atlas on \(B\) and a fixed finite atlas on \(F\), shrinking its charts
once if necessary. These charts give local ellipticity, curvature bounds
and a fixed-radius ball-volume lower bound independent of \(m\). The
same argument provides uniform harmonic-coordinate bounds if those are
desired, but harmonic coordinates are not needed for (17).

A connected compact fundamental cell for \(B_\infty\) has bounded
diameter. The cyclic chain of \(m\) cells gives
\(\operatorname{diam}(B_m)\le C m\). The circle-valued function \(z\)
modulo \(m\) has uniformly bounded differential and is onto the circle
of length \(m\), giving the reverse bound. Projection to \(B_m\) is
distance nonincreasing; a horizontal path at fixed fiber coordinate and
then a uniformly bounded fiber path prove the corresponding statements
for \(G_m\). Surjectivity also gives \(\operatorname{osc}v_m=2a\).
Equation (15) now gives the lower bound in (17).

For completeness the integral rate in (18) is exact in order, not only
an upper bound. For \(m\ge5\), deck summation gives \[
 \sum_{r=0}^{m-1}\cos^2(\theta+2\pi r/m)=\frac m2,
 \quad
 \sum_{r=0}^{m-1}\sin^4(\theta+2\pi r/m)=\frac{3m}{8}.
 \tag{20}
\] Write \(I_\eta=\int_B|\eta|^4dA_b>0\). It follows that \[
 \int_{B_m}(\Delta v_m)^2dA_{b_m}
       =\frac{a^2(2\pi)^4}{2m^3}I_\eta,
 \qquad
 \int_{B_m}|dv_m|^4dA_{b_m}
       =\frac{3a^4(2\pi)^4}{8m^3}I_\eta.
 \tag{21}
\] Insert these into (9), using the positive lower and finite upper
bounds of \(\Xi\) on \([-a,a]\), to prove (18). \(\square\)

This construction is an actual nonlinear continuation of the low channel
\(q(y)f(x)\); it is not a new smooth-topology supply theorem. It
sharpens the earlier shrinking-amplitude test by eliminating the
spurious quadratic potential of a linear metric interpolation. Its
residual has no fixed amplitude floor because every constant parameter
gives an exact Einstein product. All normal acceleration nevertheless
remains visible in (9).

\hypertarget{a-precise-fixed-distance-statement-and-its-scope}{%
\subsection{4. A precise fixed-distance statement, and its
scope}\label{a-precise-fixed-distance-statement-and-its-scope}}

Let \(\mathscr P_m\) be the metrics on \(B_m\times F\) obtained by
pulling back a hyperbolic product \(\gamma'+h_*\) through any smooth
fiber-preserving diffeomorphism \(\Phi(x,y)=(\varphi(x),\phi_x(y))\),
with orientation-preserving fiber maps. Neither the product parameters
nor the diffeomorphisms are fixed. Set \[
 d_0(G_m,\mathscr P_m)
  =\inf_{H\in\mathscr P_m}\|G_m^{-1}H-I\|_{\infty,\mathrm{op}},
 \quad
 D_T=d_T([h_{-a}],[h_a])>0,
 \tag{22}
\] where \(d_T=\tfrac12\inf\log K\) is Teichmüller distance.

\textbf{Theorem C.} For every \(m\), \[
                       d_0(G_m,\mathscr P_m)\ge\tanh D_T>0.
 \tag{23}
\] \textbf{Proof.} If \(\varepsilon=\|G_m^{-1}H-I\|_\infty\ge1\), the
claim is automatic. Otherwise
\((1-\varepsilon)G_m\le H\le(1+\varepsilon)G_m\). Restricting to a
vertical tangent space gives the same comparison between \(h_{v_m(x)}\)
and \(\phi_x^*h_*\), even when \(H\) has cross terms. The mapping class
of \(\phi_x\) is constant on connected \(B_m\), so all the latter fiber
metrics represent a single marked Teichmüller point \(t_*\). The
identity on a fiber has dilatation at most
\(\sqrt{(1+\varepsilon)/(1-\varepsilon)}\), whence \[
 d_T([h_{v_m(x)}],t_*)
             \le\frac14\log\frac{1+\varepsilon}{1-\varepsilon}.
 \tag{24}
\] Use fibers where \(v_m=\pm a\) and the triangle inequality. Then
\(D_T\le\tfrac12\log((1+\varepsilon)/(1-\varepsilon)) =\operatorname{arctanh}\varepsilon\),
proving (23). \(\square\)

Thus no estimate
\(d_0(G,\mathscr P)\le\rho(\|\operatorname{Ric}_G+G\|)\) with
\(\rho(t)\to0\) can be uniform over these cyclic covers, even for fixed
arbitrarily small but nonzero moduli amplitude, uniform local geometry,
and unnormalized \(L^2\) residual tending to zero. The omitted global
information is the slow coefficient, not a local curvature spike.

Theorem C does \textbf{not} replace \(\mathscr P_m\) by the set of all
Einstein metrics or permit arbitrary non-fiber-preserving
diffeomorphisms. Such a replacement would need a theorem showing that
every sufficiently close Einstein metric has the required global
structure. Round 18 proves local completeness for each fixed background;
I33 gives a shrinking uniform neighborhood. Neither supplies that
missing fixed-amplitude coverage. The manifolds \(M_m\) have the
explicit exact output \(b_m+h_0\), so these examples do not refute any
existence conclusion of Yau26.

\hypertarget{what-has-been-retained-rather-than-projected-away}{%
\subsection{5. What has been retained rather than projected
away}\label{what-has-been-retained-rather-than-projected-away}}

The scalar equation (6) is exact for the actual geometric family. Its
normal term is retained in (5), contributes the nonnegative \(C_N\) term
in (7), and can be nonzero. Solving \(\Delta v=0\) alone is not asserted
to solve an arbitrary full tensor perturbation problem. Within this
family, it makes \(v\) constant and then all the written vertical and
mixed equations vanish; the remaining base equation is precisely
\(K_\gamma=-1\).

This supplies a complete nonlinear energy for a genuine slow family and
a sharp global-size obstruction to uniform adapted repair. It is
compatible with the first-line obstruction/existence distinction: local
smallness and actual near-saturation may coexist with macroscopic moduli
oscillation, although the same smooth topology has an exact output. The
companion gauge calculation identifies the actual harmonic classes and
their curvature deficits; those are additional verified information, not
a presumption that every slow-moduli problem is controlled by them.

\hypertarget{reproducible-finite-checks}{%
\subsection{6. Reproducible finite
checks}\label{reproducible-finite-checks}}

\texttt{check\_nonlinear\_slow\_moduli.py} recomputes a four-dimensional
Ricci matrix from Christoffel symbols for a two-base-variable
area-preserving family, checks the normal-acceleration term, arclength
pairing and completed-square coefficients, verifies the
harmonic-coordinate cancellation and the cyclic averages in (20), and
includes negative controls for dropping acceleration or replacing the
scalar Laplacian by a Hessian norm. The local coordinate test is not a
closed hyperbolic family; the global family is supplied by the
horizontal-lift argument above. These finite checks do not certify the
Teichmüller theory or any unrestricted Einstein classification.

\end{document}
